\chapter*{\uppercase{Lampiran}}
\addcontentsline{toc}{chapter}{Lampiran}
\label{chap:Lampiran}

\section*{A. Surat Pernyataan Proyek Perancangan / Capstone}
\label{sec:Lembar_Persetujuan}

Kami yang bertanda tangan di bawah ini:

%====================================================
\begin{longtable}{|>{\centering\arraybackslash}c|
                  >{\centering\arraybackslash}p{5cm}|
                  >{\centering\arraybackslash}c|
                  >{\centering\arraybackslash}p{3.5cm}|
                  >{\centering\arraybackslash}p{2.5cm}|}
\hline
\textbf{No} &
\textbf{Nama Mahasiswa} &
\textbf{NIM} &
\textbf{Program Studi} &
\textbf{Peran} \\ \hline

1 & & & & \\[0.6cm] \hline
2 & & & & \\[0.6cm] \hline
3 & & & & \\[0.6cm] \hline
4 & & & & \\[0.6cm] \hline

\end{longtable}

Menyatakan bahwa Proyek Perancangan / Capstone yang diajukan :
\vspace{0.3cm}

\textbf{Judul Capstone} : \\[0.1cm]
\hrulefill \\[0.4cm]

\textbf{Tema Capstone} : \\[0.1cm]
\hrulefill \\[0.5cm]



%====================================================
\begin{longtable}{cp{13cm}}
1 & Merupakan karya asli tim dan bebas plagiarisme \\[0.2cm]
2 & Mengandung unsur engineering design untuk menyelesaikan masalah nyata \\[0.2cm]
3 & Mempertimbangkan aspek teknis, ekonomi, keselamatan, lingkungan, sosial, dan etika \\[0.2cm]
4 & Memenuhi kategori complex engineering problem sesuai CPL/PLO \\[0.2cm]
5 & Dikerjakan secara kolaboratif dengan pembagian tugas yang jelas \\[0.2cm]
6 & Menggunakan metode analisis, eksperimen, atau simulasi yang relevan \\[0.2cm]
7 & Menghasilkan prototipe, model, sistem, atau rekomendasi implementatif \\[0.2cm]
8 & Seluruh data dan referensi digunakan secara bertanggung jawab \\[0.2cm]
9 & Laporan dan presentasi disusun secara mandiri oleh tim \\[0.2cm]
10 & Bersedia dijadikan evidensi akreditasi IABEE \\

\end{longtable}

\vspace{0.3cm}

\textbf{Tim Proyek Perancangan / Capstone}

\begin{longtable}{|p{7cm}|p{7cm}|}
\hline
\centering
\textbf{Nama Mahasiswa} & 
\centering
\textbf{Tanda Tangan} \\ 
\hline
& \\[0.5cm] \hline
& \\[0.5cm] \hline
& \\[0.5cm] \hline
& \\[0.5cm] \hline
\end{longtable}
\vspace{0.1cm}
%====================================================
\section*{B. Pemetaan Capaian Pembelajaran}

\begin{longtable}{|c|p{10cm}|c|}
\hline
\centering
\textbf{CPL} & 
\centering 
\textbf{Deskripsi Kriteria} & \textbf{Konfirmasi Pembimbing} \\ \hline
CPL-1 & Penerapan matematika, sains, dan engineering & $\square$ \\ \hline
CPL-2 & Analisis masalah kompleks & $\square$ \\ \hline
CPL-3 & Desain solusi engineering & $\square$ \\ \hline
CPL-4 & Investigasi / eksperimen & $\square$ \\ \hline
CPL-5 & Penggunaan tools modern & $\square$ \\ \hline
CPL-6 & Kerja tim & $\square$ \\ \hline
CPL-7 & Komunikasi efektif & $\square$ \\ \hline
CPL-8 & Etika profesional & $\square$ \\ \hline
CPL-9 & Dampak sosial dan lingkungan & $\square$ \\ \hline
CPL-10 & Manajemen proyek & $\square$ \\ \hline
\end{longtable}

\section*{C. Pembuktian Konsep Prototipe}
\textbf{Demo} : Lulus / Tidak Lulus


%Tempat, Tanggal : \hrulefill

\begin{flushright}
Tempat, Tanggal\\[0.2cm]
\vspace{0.2cm}
Mengetahui, \\[0.2cm]
\textbf{Dosen Pembimbing} \\[2cm]

Nama : .................................... \\[0.2cm]
NIDN/NIP : ....................................
\end{flushright}